\documentclass[10pt,reqno]{amsart}

\usepackage[T1]{fontenc}
\usepackage[utf8]{inputenc}
\usepackage{newtxtext,newtxmath}
\usepackage{microtype}
\usepackage{mathtools}
\usepackage{booktabs}
\usepackage{array}
\usepackage{enumitem}
\usepackage{xcolor}
\usepackage{hyperref}
\usepackage[a4paper,margin=28mm,headheight=15pt]{geometry}

\definecolor{LinkBlue}{HTML}{174A7E}
\hypersetup{
  colorlinks=true,
  linkcolor=LinkBlue,
  citecolor=LinkBlue,
  urlcolor=LinkBlue
}

\setlength{\parskip}{2pt}
\setlist[itemize]{leftmargin=6mm,itemsep=1.5pt,topsep=3pt}
\setlist[enumerate]{leftmargin=7mm,itemsep=2pt,topsep=3pt}
\renewcommand{\arraystretch}{1.15}
\emergencystretch=2em

\newtheorem{theorem}{Theorem}[section]
\newtheorem{proposition}[theorem]{Proposition}
\newtheorem{lemma}[theorem]{Lemma}
\newtheorem{corollary}[theorem]{Corollary}
\theoremstyle{definition}
\newtheorem{definition}[theorem]{Definition}
\newtheorem{remark}[theorem]{Remark}

\title[Distinct Hilbert-Scheme Components Sharing an Unlabelled Gr\"obner Fan]
{Distinct Hilbert-Scheme Components Sharing a Complete Unlabelled Gr\"obner Fan}
\author{Argus Open Research Project}
\date{September 8, 2026}
\subjclass[2020]{Primary 14C05; Secondary 13P10, 52B20}
\begin{document}
\begin{abstract}
Problem~15 (AIMPL~1.35) asks whether the Gr\"obner fan is a discrete invariant capable of distinguishing distinct irreducible components of the Hilbert scheme $\mathrm{Hilb}^d(\mathbb{A}^n)$. Under the fixed-coordinate unlabelled geometric Gr\"obner fan convention of Fukuda, Jensen, and Thomas, we construct conjugate frame quotients of length $d=4803$ in $\mathbb{A}^{20}_{\mathbb{C}}$ defined over $\mathbb{Q}(\sqrt{-7})$ associated with nonisomorphic rigid projective surfaces without real structure. We prove that they lie on distinct entire irreducible components of $\mathrm{Hilb}^{4803}(\mathbb{A}^{20}_{\mathbb{C}})$ that admit simultaneous component-exclusive nonempty dense open subsets on which every ideal has the exact same complete fixed-coordinate unlabelled Gr\"obner fan.
\end{abstract}
\maketitle
\section{Main Theorem}
Let $S=\mathbb{C}[x_1,\dots,x_n]$. For an ideal $I\subset S$, the geometric Gr\"obner fan $\mathrm{GF}(I)$ in $\mathbb{R}_{\ge 0}^n$ is the polyhedral fan of weight cones $C_\omega(I)=\overline{\{\omega'\in\mathbb{R}_{>0}^n : \mathrm{in}_{\omega'}(I)=\mathrm{in}_\omega(I)\}}$.
\begin{theorem}
For $n=20$ and $d=4803$, there exist two distinct irreducible components $Z^+, Z^-\subset\mathrm{Hilb}^{4803}(\mathbb{A}^{20}_{\mathbb{C}})$ and nonempty Zariski-open subsets $U^+\subset Z^+\setminus\bigcup_{W\neq Z^+} W$ and $U^-\subset Z^-\setminus\bigcup_{W\neq Z^-} W$ such that for every $[I^+]\in U^+$ and $[I^-]\in U^-$, the complete fixed-coordinate unlabelled Gr\"obner fans coincide:
\[
  \mathrm{GF}(I^+) \;=\; \mathrm{GF}(I^-).
\]
\end{theorem}
\begin{proof}[Proof summary]
Let $X^+, X^-$ be a complex-conjugate pair of rigid surfaces of general type defined over $\mathbb{Q}(\sqrt{-7})$ admitting no real structure and no birational isomorphism $X^+\dashrightarrow X^-$. Taking affine cones and a generic linear frame truncation yields conjugate Artinian quotients $S/I_0^+$ and $S/I_0^-$ of colength $d=4803$ in $n=20$ variables. A rational ambient-open reconstruction map recovers the rigid surface isomorphism class from any deformation in the same component, separating every minimal component $Z^+$ containing $[I_0^+]$ from its complex conjugate $Z^-=\overline{Z^+}$. On the other hand, complex conjugation of polynomial coefficients preserves initial monomial ideals $\mathrm{in}_\omega(\bar{I})=\mathrm{in}_\omega(I)$ for every weight vector $\omega\in\mathbb{R}^n$, and a common degree-six nonvanishing minor condition cuts out simultaneous component-exclusive dense opens $U^\pm\subset Z^\pm$ with constant complete Gr\"obner fan.
\end{proof}
\end{document}
